\section{Appendix A. Mathematical Derivations}

\subsection{A.1 Truncated SVD}

For $W \in \mathbb{R}^{m \times n}$, the SVD is:

\begin{equation}
W = U \Sigma V^\top
\end{equation}

The best rank-$k$ approximation in Frobenius norm (Eckart-Young):

\begin{equation}
W_k = U_{:, :k} \Sigma_{:k, :k} V_{:k, :}^\top
\end{equation}

with error:

\begin{equation}
\|W - W_k\|_F = \sqrt{\sum_{i > k} \sigma_i^2}
\end{equation}

\subsection{A.2 Factorization}

We split (A.2) into two factors:

\begin{equation}
\text{down} = V_{:k, :}^\top \in \mathbb{R}^{k \times n}
\end{equation}

\begin{equation}
\text{up} = U_{:, :k} \Sigma_{:k, :k} \in \mathbb{R}^{m \times k}
\end{equation}

Then:

\begin{equation}
\text{up} \cdot \text{down} = U_{:, :k} \Sigma_{:k, :k} V_{:k, :}^\top = W_k
\end{equation}

\subsection{A.3 Parameter Count}

For each factorized projection:

\begin{equation}
\text{params} = k \cdot n + m \cdot k = k(m + n)
\end{equation}

vs baseline $m \cdot n$. Reduction factor:

\begin{equation}
\frac{m \cdot n}{k(m + n)}
\end{equation}

For square $m = n = d$:

\begin{equation}
\frac{d^2}{2kd} = \frac{d}{2k}
\end{equation}

\subsection{A.4 LoRA Composition}

With LoRA attached to the up-projection:

\begin{equation}
\text{out} = (W_{up} + \Delta W_{up}) \cdot \text{down}(x)
\end{equation}

where:

\begin{equation}
\Delta W_{up} = \frac{\alpha}{r_L} B A
\end{equation}

So the effective weight is:

\begin{equation}
W_\text{eff} = (W_{up} + \frac{\alpha}{r_L} B A) \cdot W_\text{down}
\end{equation}

\subsection{A.5 GQA Parameter Reduction}

For GQA with $H_{kv} < H$:

\begin{equation}
\text{K params baseline} = d \cdot (H_{kv} \cdot d_h) = d^2 \cdot (H_{kv}/H)
\end{equation}

\begin{equation}
\text{K params ours} = k(d + H_{kv} d_h)
\end{equation}

For $H_{kv} = H/4$ and $k = d/6$:

\begin{equation}
\text{reduction} = \frac{d^2 (H_{kv}/H)}{k(d + H_{kv} d_h)} \approx \frac{d}{k} \cdot \frac{1}{1 + H_{kv} d_h / d}
\end{equation}

\subsection{A.6 Effect of Rank on Error}

The Frobenius error after rank-$k$ truncation is:

\begin{equation}
\text{err}(k) = \sqrt{\sum_{i > k} \sigma_i^2} / \|W\|_F
\end{equation}

For transformer attention weights, this decays rapidly with $k$.

\subsection{A.7 Symmetric vs Asymmetric Split}

Our implementation uses asymmetric split ($\Sigma$ entirely in the up-projection):

\begin{equation}
\text{down} = V_k^\top, \quad \text{up} = U_k \Sigma_k
\end{equation}

An alternative symmetric split:

\begin{equation}
\text{down} = \Sigma_k^{1/2} V_k^\top, \quad \text{up} = U_k \Sigma_k^{1/2}
\end{equation}

Both give the same $UV$, but asymmetric has:
\begin{itemize}
\item Slightly better numerical behavior in FP16.
\item All scaling concentrated in one factor.

\end{itemize}
